\chapter{Loss Function Definitions and Gradients}

\label{chap:loss-definitions}

This appendix expands the loss-family definitions from Chapter~3 \S3.3 with explicit closed-form expressions and derivatives with respect to the prediction $\hat y$. Every formula in Sections A.1--A.3.2 matches the implementation in \texttt{Model\_Train/losses.py} on the current \texttt{main} branch. Section~A.3.3 records the conceptual form used to interpret Phase~2 branch outputs; its exact closed form lives on the \texttt{phase2.2-fix} branch. Gradients are stated pointwise; the effective training gradient at each optimiser step is the batch average (or median for \texttt{MedSE}) of the pointwise gradient.

Throughout this appendix, $y$ is the scalar realised return, $\hat y$ is the scalar prediction, and $e = y - \hat y$ is the residual. Where a batch-normalisation factor appears, the batch-level mean is written $\mathbb{E}_{\mathcal{B}}[\cdot]$ and the small regulariser $\epsilon = 10^{-8}$ is included to avoid division by zero.

% ============================================================
\section{Regression Losses}

\label{sec:regression-losses}

\subsection{MSE}

\label{subsec:mse}

\begin{equation}
L_{\mathrm{MSE}}(y, \hat y) = (y - \hat y)^2,
\qquad
\frac{\partial L_{\mathrm{MSE}}}{\partial \hat y} = -2\,(y - \hat y).
\end{equation}

Batch reduction: arithmetic mean.

\subsection{MedSE}

\label{subsec:medse}

\begin{equation}
L_{\mathrm{MedSE}}(y, \hat y) = (y - \hat y)^2,
\qquad
\text{batch reduction} = \mathrm{median}.
\end{equation}

Because the median is non-decomposable, the gradient at each observation is zero at every observation \emph{except} those that determine the batch median. In PyTorch the autograd routes gradient only through the median-ranked observation; this is why MedSE training is inherently noisier and why \texttt{reduction="median"} is declared explicitly in \texttt{medse\_loss}.

\subsection{Huber (as magnitude backbone)}

\label{subsec:huber}

Used inside every hybrid loss, with $\delta = 0.01$ hard-coded in \texttt{\_huber\_term}:

\begin{equation}
H_\delta(e) =
\begin{cases}
\tfrac12 e^2 & |e| \le \delta, \\
\delta\, (|e| - \tfrac12 \delta) & |e| > \delta,
\end{cases}
\qquad
\frac{\partial H_\delta}{\partial \hat y}
=
\begin{cases}
-e & |e| \le \delta, \\
-\delta \cdot \operatorname{sgn}(e) & |e| > \delta.
\end{cases}
\end{equation}

The gradient is continuous in $e$ (the one-sided derivatives agree at both boundaries $|e| = \delta$) and bounded by $\delta$ in absolute value.

% ============================================================
\section{Directional Losses}

\label{sec:directional-losses}

\subsection{MADL (differentiable)}

\label{subsec:madl}

With $a = 25$ fixed:

\begin{equation}
L_{\mathrm{MADL}}(y, \hat y) = -\tanh(a \cdot y \cdot \hat y) \cdot |y|,
\end{equation}
\begin{equation}
\frac{\partial L_{\mathrm{MADL}}}{\partial \hat y} = -a \cdot y \cdot |y| \cdot \operatorname{sech}^2(a \cdot y \cdot \hat y).
\end{equation}

The gradient peaks near $a \cdot y \cdot \hat y = 0$ and decays exponentially as $|a \cdot y \cdot \hat y|$ grows. For realistically small $|y|$ (monthly returns typically a few per cent), the factor $y \cdot |y| = y^2 \cdot \operatorname{sgn}(y)$ is of order $10^{-4}$, which is the ``weak-gradient-near-zero'' behaviour described in Chapter~2 \S2.4.

\subsection{GMADL}

\label{subsec:gmadl}

With $a = 100$ and $b = 2$ fixed:

\begin{equation}
L_{\mathrm{GMADL}}(y, \hat y) = -\big[\sigma(a \cdot y \cdot \hat y) - \tfrac12\big] \cdot |y|^b,
\end{equation}
\begin{equation}
\frac{\partial L_{\mathrm{GMADL}}}{\partial \hat y}
= -a \cdot y \cdot |y|^b \cdot \sigma(a \cdot y \cdot \hat y) \cdot \big[1 - \sigma(a \cdot y \cdot \hat y)\big].
\end{equation}

The $|y|^b = y^2$ factor provides stronger magnitude weighting than MADL, but the sigmoid derivative $\sigma(1-\sigma)$ saturates at $|a y \hat y| \gg 1$. Together these explain why GMADL produces larger absolute prediction values (to push $|a y \hat y|$ into the saturation region) while preserving directional ranking --- and why average R\textsuperscript{2} values diverge under this loss (Chapter~5 \S5.2).

\subsection{IMADL (rebalanced)}

\label{subsec:imadl}

With $a = 100$, $b = 2$, $\lambda_{\mathrm{dir}} = \lambda_{\mathrm{mag}} = 1$ fixed:

\begin{equation}
D(y, \hat y) = \big[1 - \sigma(a\, y\, \hat y)\big] \cdot \frac{|y|^b}{\mathbb{E}_{\mathcal{B}}[|y|^b] + \epsilon},
\end{equation}
\begin{equation}
L_{\mathrm{IMADL}}(y, \hat y) = \lambda_{\mathrm{dir}} \cdot D(y, \hat y) + \lambda_{\mathrm{mag}} \cdot (y - \hat y)^2,
\end{equation}
\begin{multline}
\frac{\partial L_{\mathrm{IMADL}}}{\partial \hat y}
= -\lambda_{\mathrm{dir}} \cdot a\, y \cdot \sigma(a\, y\, \hat y) \cdot \big[1 - \sigma(a\, y\, \hat y)\big] \cdot \frac{|y|^b}{\mathbb{E}_{\mathcal{B}}[|y|^b] + \epsilon} \\
- 2\,\lambda_{\mathrm{mag}}\,(y - \hat y).
\end{multline}

The batch-normalisation factor $\mathbb{E}_{\mathcal{B}}[|y|^b]$ decouples the scale of the directional gradient from the particular batch composition; it is treated as a constant during the backward pass through that observation (for gradients with respect to $\hat y$, the normalisation denominator is constant, matching the design intent).

% ============================================================
\section{Hybrid Losses}

\label{sec:hybrid-losses}

\subsection{Additive (A-series)}

\label{subsec:additive}

With Huber $\delta = 0.01$ fixed and $(\lambda_{\mathrm{dir}}, \lambda_{\mathrm{hub}})$ varied per variant (see Chapter~3 \S3.3.3):

\begin{equation}
L_{\mathrm{add}}(y, \hat y) = \lambda_{\mathrm{dir}} \cdot D(y, \hat y) + \lambda_{\mathrm{hub}} \cdot H_\delta(y - \hat y),
\end{equation}
\begin{equation}
\frac{\partial L_{\mathrm{add}}}{\partial \hat y}
= \lambda_{\mathrm{dir}} \cdot \frac{\partial D}{\partial \hat y}
+ \lambda_{\mathrm{hub}} \cdot \frac{\partial H_\delta}{\partial \hat y}.
\end{equation}

Both components contribute additively to the gradient; the relative weighting is $(\lambda_{\mathrm{dir}}, \lambda_{\mathrm{hub}})$. The A-variants scan this 2-dimensional hyperparameter.

\subsection{Multiplicative (M-series)}

\label{subsec:multiplicative}

\begin{equation}
L_{\mathrm{mul}}(y, \hat y) = \big(1 + \lambda_{\mathrm{dir}} \cdot D(y, \hat y)\big) \cdot H_\delta(y - \hat y),
\end{equation}
\begin{equation}
\frac{\partial L_{\mathrm{mul}}}{\partial \hat y}
= \lambda_{\mathrm{dir}} \cdot \frac{\partial D}{\partial \hat y} \cdot H_\delta(y - \hat y)
+ \big(1 + \lambda_{\mathrm{dir}} \cdot D(y, \hat y)\big) \cdot \frac{\partial H_\delta}{\partial \hat y}.
\end{equation}

The multiplicative form couples the two components: the Huber gradient is \emph{gated} by the directional factor $1 + \lambda_{\mathrm{dir}} D$. When the direction is predicted correctly, $D \to 0$ and the gradient reduces to $\partial H_\delta / \partial \hat y$. When the direction is predicted wrongly, the sigmoid penalty $\sigma(ay\hat y)$ approaches 0 so $[1-\sigma(\cdot)]$ approaches 1, and the Huber gradient is amplified by a factor of $1 + \lambda_{\mathrm{dir}} \cdot D$, where $D$ is proportional to the batch-normalised magnitude weight $|y|^b / (\mathbb{E}_{\mathcal{B}}[|y|^b]+\epsilon)$.

\subsection{M2-robust \texorpdfstring{$\gamma$}{gamma} family}

\label{subsec:m2-robust}

Conceptually, the M2-robust loss replaces the fixed-threshold Huber backbone with a $\gamma$-controlled saturating magnitude term $H^{\gamma}(e)$ and combines it with the same directional-gating structure:

\begin{equation}
L_{\mathrm{M2\text{-}robust},\gamma}(y, \hat y) = \big(1 + \lambda_{\mathrm{dir}} \cdot D(y, \hat y)\big) \cdot H^{\gamma}(y - \hat y).
\end{equation}

The robustness schedule $H^{\gamma}(e)$ is a smooth monotone function that approaches $\tfrac12 e^2$ for small $|e|$ and saturates above a threshold controlled by $\gamma$. Small $\gamma$ flattens the loss surface and prioritises the directional component; large $\gamma$ recovers a nearly unbounded quadratic behaviour. The $\gamma$ refinement of Chapter~5 \S5.4 scans $\gamma \in \{0.03, 0.05, 0.07, 0.10, 0.15\}$. The exact closed form used in the Phase~2.2 runner lives on the \texttt{phase2.2-fix} branch together with the runner scripts; the empirically selected setting within the reported 24-month, 3-seed protocol is $\gamma = 0.07$ (see Chapter~5 \S5.4 and Table~5.3).

% ============================================================
\section{Notes on Practical Implementation}

\label{sec:practical-notes}

Two implementation points are worth repeating for completeness, both documented in \texttt{Model\_Train/losses.py}:

\begin{enumerate}
    \item \textbf{Reduction.} Every loss function uses \texttt{reduction="mean"} by default, except MedSE which uses \texttt{reduction="median"}. Passing \texttt{reduction="none"} returns the pointwise loss vector, which the tests in \texttt{tests/} consume.
    \item \textbf{Numerical guards.} IMADL and the hybrid variants use the batch-mean normalisation $\mathbb{E}_{\mathcal{B}}[|y|^b] + \epsilon$ with $\epsilon = 10^{-8}$; the Huber term has no special numerical guard because it is smooth and bounded.
\end{enumerate}
